% =====================================================================
%  DAFTAR PUSTAKA
%  Gaya: Harvard Anglia Ruskin University (ARU)
%  Format manual — tanpa BibTeX
%  Urutan: alfabet nama belakang penulis pertama
% =====================================================================

\chapter*{Daftar Pustaka}
\addcontentsline{toc}{chapter}{Daftar Pustaka}

\begingroup
\setlength{\parindent}{-0.5cm}
\setlength{\leftskip}{0.5cm}
\setlength{\parskip}{6pt}

% ── FORMAT HARVARD ARU ────────────────────────────────────────────────
%
% JURNAL
%   Nama Belakang, Inisial. (Tahun) 'Judul artikel',
%   \textit{Nama Jurnal}, Volume(Nomor), hlm. XX--XX.
%   doi: xxxxx [Diakses: DD Bulan YYYY].
%
% BUKU
%   Nama Belakang, Inisial. (Tahun) \textit{Judul Buku}.
%   Edisi ke-X. Kota: Penerbit.
%
% BAB DALAM BUKU YANG DIEDIT
%   Nama Belakang, Inisial. (Tahun) 'Judul bab', dalam
%   Nama Belakang Editor, Inisial. (ed./eds.)
%   \textit{Judul Buku}. Kota: Penerbit, hlm. XX--XX.
%
% PROSIDING KONFERENSI
%   Nama Belakang, Inisial. (Tahun) 'Judul makalah',
%   \textit{Nama Konferensi}, Kota, DD--DD Bulan YYYY.
%   Kota: Penerbit, hlm. XX--XX.
%
% SITUS WEB / DOKUMEN ONLINE
%   Nama Belakang, Inisial. / Nama Organisasi (Tahun)
%   \textit{Judul Halaman}. Tersedia di: URL
%   [Diakses: DD Bulan YYYY].
%
% DUA ATAU TIGA PENULIS: pisah dengan 'dan'
%   Doe, J. dan Smith, A. (Tahun) ...
%
% EMPAT PENULIS ATAU LEBIH: tulis semua sampai penulis terakhir
%   (tidak disingkat \textit{et al.} dalam daftar pustaka —
%    \textit{et al.} hanya untuk sitasi dalam teks)
% ─────────────────────────────────────────────────────────────────────

% Contoh jurnal:
[Nama Belakang], [I]. ([Tahun]) '[Judul artikel]',
\textit{[Nama Jurnal]}, [Vol]([No]), hlm.~[XX--XX].
doi:~[xxxxx] [Diakses: DD Bulan YYYY].

% Contoh buku:
[Nama Belakang], [I]. ([Tahun]) \textit{[Judul Buku]}.
[Kota]: [Penerbit].

% Tambahkan referensi di bawah, diurutkan alfabet nama belakang penulis pertama.

\endgroup
